% !PRINTED 366
\weberpage{366}
Now we have the two complexes:

\weq{4}{x_2 z_2,\ x_3 z_3,\ \dots}
\weq{4\text{b}}{x_2 x_3,\ z_2 z_3,\ \dots{},}

and since $x_2 x_3 y_2 y_3$ are a syzygetic quadruple, the group (4 b) also contains the pair $y_2 y_3$, and consequently $y_2 y_3 z_2 z_3$ are also a syzygetic quadruple. From this it follows further that $y_2 z_2$ and $y_3 z_3$ belong to the same complex. Since the same consideration as for $x_3 y_3$ can also be carried out for the pairs $x_4 y_4, x_5 y_5, x_6 y_6$, the complexes (4), (5) can accordingly be formed completely, and they take the expression:

\weq{4}{x_1 z_1,\ x_2 z_2,\ x_3 z_3,\ x_4 z_4,\ x_5 z_5,\ x_6 z_6}
\weq{5}{y_1 z_1,\ y_2 z_2,\ y_3 z_3,\ y_4 z_4,\ y_5 z_5,\ y_6 z_6.}

Here $z_1 z_2$ form a pair of the complex (2), and $z_3, z_4, z_5, z_6$, which form an azygetic system, all occur in the complex (3), in which no two are paired.

Since, as shown above, we can derive from the arbitrarily assumed complex (1), after the type of (2), (3), (4), (5), all complexes that exist at all, the theorem results:

6. \gesperrt{Any two complexes have either a syzygetic quadruple or an azygetic system of six double tangents in common.}

Two complexes that have a syzygetic quadruple in common we shall call a \gesperrt{syzygetic complex pair}. To every such pair there is a third complex which contains the same quadruple. Three such complexes we call a \gesperrt{syzygetic complex triple} [e.\,g. (1), (2), (3)].

Likewise we call two complexes of the second kind, that is, such as have six azygetic elements in common, an \gesperrt{azygetic complex pair}.

Every azygetic complex pair is likewise completed to a triple by a determinate complex which forms an azygetic pair with each of the two complexes [as (1), (4), (5)]. Such a one we call an \gesperrt{azygetic complex triple}.

In a syzygetic triple all 28 double tangents occur, in an azygetic one only 18.
